\subsection{Homogeneous prime supports of projective seminorms}
\label{subsec:projective-support}

The support is the actual prime kernel of a seminorm. We first construct it
for arbitrary internally graded commutative rings, and then specialize to the
normalized projective model. The target is Mathlib's projective spectrum of
relevant homogeneous prime ideals with its Zariski topology; no identification
of a scheme's analytification or structure sheaf is assumed.

\begin{lemma}[Homogeneous kernels of graded seminorms]
\label{spectrum:graded-support}
\lean{EnvelopeProperty.Berkovich.Spectrum.IsGraded,
EnvelopeProperty.Berkovich.Spectrum.supp_isHomogeneous,
EnvelopeProperty.Berkovich.Spectrum.homogeneousSupport,
EnvelopeProperty.Berkovich.Spectrum.homogeneousSupport_toIdeal,
EnvelopeProperty.Berkovich.Spectrum.primePoint_val_eq_one_of_not_mem,
EnvelopeProperty.Berkovich.Spectrum.primePoint_val_le_one,
EnvelopeProperty.Berkovich.Spectrum.primePoint_isGraded}
\leanok
\uses{spectrum:support-surjective}
Let $A=\bigoplus_{i\in\Gamma} A_i$ be an internally graded commutative ring,
where $\Gamma$ is an additive monoid. Call a real multiplicative ultrametric
seminorm $x$ graded if $|a_i|_x\le |a|_x$ for every homogeneous projection
$a_i$ of every $a\in A$. Its prime kernel is a homogeneous ideal.
For any homogeneous prime $\mathfrak p$, the seminorm pulled back from the
trivial seminorm on $A/\mathfrak p$ is graded and has kernel $\mathfrak p$.
\end{lemma}
\begin{proof}
If $|a|_x=0$, the grading inequalities give $|a_i|_x=0$ for every $i$.
Thus the kernel is closed under homogeneous projections, exactly the
homogeneous-ideal condition; it is prime by multiplicativity.
The trivial prime-quotient seminorm takes value zero precisely on
$\mathfrak p$, and value one elsewhere. If $a\in\mathfrak p$, homogeneity
places each $a_i$ in $\mathfrak p$, so both values are zero.
If $a\notin\mathfrak p$, its value is one and every projection has value
at most one. This proves all grading inequalities.
\end{proof}

\begin{theorem}[Supports for arbitrary relevant graded spectra]
\label{spectrum:relevant-graded-support}
\lean{EnvelopeProperty.Berkovich.Spectrum.projectiveSupport,
EnvelopeProperty.Berkovich.Spectrum.projectiveSupport_toIdeal,
EnvelopeProperty.Berkovich.RelevantGradedSpectrum,
EnvelopeProperty.Berkovich.RelevantGradedSpectrum.support,
EnvelopeProperty.Berkovich.RelevantGradedSpectrum.support_toIdeal,
EnvelopeProperty.Berkovich.RelevantGradedSpectrum.mem_support,
EnvelopeProperty.Berkovich.RelevantGradedSpectrum.continuous_support,
EnvelopeProperty.Berkovich.RelevantGradedSpectrum.primePoint,
EnvelopeProperty.Berkovich.RelevantGradedSpectrum.support_primePoint,
EnvelopeProperty.Berkovich.RelevantGradedSpectrum.support_surjective}
\leanok
\uses{spectrum:graded-support,spectrum:support}
For an $\mathbb N$-graded commutative ring $A$, consider the subspace of
graded real seminorms whose kernels do not contain the irrelevant ideal
$A_+$. Sending a point to its prime kernel gives a continuous surjection
onto the projective spectrum $\operatorname{Proj}(A)$.
\end{theorem}
\begin{proof}
The preceding lemma bundles each kernel as a homogeneous prime. The
explicit irrelevant-avoidance condition makes it a projective-spectrum point.
The inverse image of the Zariski zero locus of a set $E\subseteq A$ is
$\bigcap_{a\in E}\{x:|a|_x=0\}$, a closed set by continuous evaluation.
Every Zariski closed set is such a zero locus, proving continuity.
Given a relevant homogeneous prime, the trivial prime-quotient seminorm
is graded by the preceding lemma, has precisely that kernel, and therefore
satisfies the relevance condition. This gives surjectivity.
\end{proof}

\begin{theorem}[Projective-model support]
\label{projective:support}
\lean{EnvelopeProperty.ProjBerk.proj_eq_homogeneousComponent,
EnvelopeProperty.ProjBerk.irrelevant_eq_idealOfVars,
EnvelopeProperty.ProjBerk.toRelevantGradedSpectrum,
EnvelopeProperty.ProjBerk.continuous_toRelevantGradedSpectrum,
EnvelopeProperty.ProjBerk.support,
EnvelopeProperty.ProjBerk.support_toIdeal,
EnvelopeProperty.ProjBerk.mem_support,
EnvelopeProperty.ProjBerk.continuous_support,
EnvelopeProperty.ProjBerk.preimage_support_zeroLocus,
EnvelopeProperty.ProjBerk.exists_X_not_mem,
EnvelopeProperty.ProjBerk.homogeneousPrimePoint,
EnvelopeProperty.ProjBerk.support_homogeneousPrimePoint,
EnvelopeProperty.ProjBerk.support_surjective}
\leanok
\uses{spectrum:relevant-graded-support,charts:forward-chart,berk:support-surjective,
closed-charts:vanishing}
For a trivially valued field $k$, the normalized projective seminorm model
has a continuous surjective support map to the projective spectrum of
$k[X_0,\ldots,X_n]$ with its total-degree grading. Its homogeneous prime
ideal is exactly the seminorm kernel, and the preimage of a projective
Zariski zero locus is its analytic evaluation zero locus.
\end{theorem}
\begin{proof}
Gradedness of a normalized projective point gives the homogeneous-component
inequalities, so its kernel is homogeneous. Some coordinate has value one
by normalization. That coordinate is a positive-degree homogeneous element
and hence lies in the irrelevant ideal, but not in the kernel. This excludes
the irrelevant ideal and defines the projective-spectrum point. The preceding
generic support construction and the continuous inclusion of the projective
model prove continuity; the zero-locus formula follows from the kernel definition.
For surjectivity, let $\mathfrak p$ be a relevant homogeneous prime. At least
one coordinate is outside $\mathfrak p$, because the coordinate ideal contains
every positive-degree homogeneous polynomial and therefore the irrelevant ideal.
The trivial prime-quotient seminorm is graded, takes coordinate values zero
or one, and takes value one on a coordinate outside $\mathfrak p$.
Consequently its maximum coordinate value is one, so it is a normalized
projective point with support $\mathfrak p$.
\end{proof}
